\chapter{晚清：洋务、甲午与新政}
\label{ch:late-qing-reform}
\begin{figure}[htbp]
  \centering
  \begin{tikzpicture}[
    x=1cm,y=1cm,
    date/.style={font=\scriptsize\bfseries, text=dynasty},
    event/.style={draw=timeline, fill=paper, rounded corners=2pt, align=center,
      inner sep=3pt, font=\scriptsize, text=ink},
    note/.style={font=\scriptsize\color{source}, align=center}
  ]
    \draw[timeline, line width=1.2pt] (0,0) -- (12,0);
    \foreach \x/\year in {0/1878,1.8/1884--85,3.6/1894--95,5.5/1898,7.3/1900--01,9.2/1905,10.6/1906,12/1908}
      {\draw[timeline] (\x,.12) -- (\x,-.12); \node[date, below] at (\x,-.18) {\year};}
    \node[event, above] at (0,.42) {新疆战事收束；\\设省讨论};
    \node[event, below] at (1.8,-.62) {中法战争、台湾\\建省与海防调整};
    \node[event, above] at (3.6,.42) {甲午战争、\\《马关条约》};
    \node[event, below] at (5.5,-.62) {戊戌变法\\与政变};
    \node[event, above] at (7.3,.42) {义和团战争、\\《辛丑条约》};
    \node[event, below] at (9.2,-.62) {日俄战争；\\废科举};
    \node[event, above] at (10.6,.42) {预备立宪、\\新军与学制改革};
    \node[event, below] at (12,-.62) {宪法大纲；\\摄政与继承};
    \node[note] at (6,-1.35)
      {设省、建军、学校、铁路、条约和立宪的制度日期，不等于地方采纳或社会影响的同一日期。};
  \end{tikzpicture}
  \caption{晚清改革与边政的编年定位（1878--1908，简表）。图把新疆、台湾、东北和沿海的边政--战争节点，同甲午、义和团、新政和立宪程序并列。主要依据《清德宗实录》《清宣统政纪》、军机处和总理衙门档、条约文本、机构记录与中外报刊；诏令、方案、编练和工程开办只能证明相应制度动作，不能证明全省覆盖或实际效果。}
  \label{fig:late-qing-reform-timeline}
\end{figure}

\begin{figure}[htbp]
  \centering
  \begin{tikzpicture}[
    x=1cm,y=1cm,
    place/.style={draw=timeline, fill=paper, rounded corners=2pt, align=center,
      inner sep=3pt, font=\small},
    court/.style={place, draw=dynasty, fill=dynasty!13},
    reform/.style={place, draw=timeline, fill=timeline!13},
    frontier/.style={place, draw=source, fill=source!10},
    note/.style={align=center, font=\scriptsize\color{source}}
  ]
    \fill[paper] (-.65,-3.9) rectangle (7.85,3.85);
    \draw[dynasty, line width=.8pt] (-.65,3.85) -- (7.85,3.85);
    \node[font=\bfseries\large, text=dynasty] at (3.6,3.42)
      {1878--1908：新省、军政与新政的省际节点};

    \node[court] (beijing) at (3.55,2.55) {北京\\部院、总理衙门\\新政与立宪诏令};
    \node[frontier] (xinjiang) at (.2,1.1) {新疆\\收复后设省\\1878--1884};
    \node[frontier] (northeast) at (6.5,1.35) {东三省\\日俄战争后\\行省体制调整};
    \node[reform] (tianjin) at (5.35,.5) {天津、直隶\\北洋军政、学堂\\与铁路节点};
    \node[reform] (shanghai) at (2.1,.25) {上海、江南\\工商、报刊、\\教育与法律讨论};
    \node[reform] (fuzhou) at (5.95,-1.65) {福州、台湾海峡\\船政、台湾建省\\电报与铁路试验};
    \node[reform] (wuhan) at (3.35,-1.3) {武昌、湖北\\新军、学堂与\\省级行政实验};
    \node[reform] (guangzhou) at (.35,-2.5) {两广、香港邻近\\商政、教育与\\海上信息网络};
    \node[frontier] (yunnan) at (-.15,-.75) {云南、越南边境\\中法战争与\\边政重组};

    \draw[dynasty, dashed] (beijing) -- node[above,sloped,note]
      {诏令、经费与任命的路径} (tianjin);
    \draw[dynasty, dashed] (beijing) -- node[above,sloped,note]
      {边政、军费与省制} (xinjiang);
    \draw[dynasty, dashed] (beijing) -- node[above,sloped,note]
      {新政与东北改制} (northeast);
    \draw[->, timeline, line width=1.1pt] (tianjin) -- node[right,sloped,note]
      {人员、报刊与铁路的多向流动} (shanghai);
    \draw[->, timeline, line width=1.1pt] (shanghai) -- node[right,sloped,note]
      {工商、教育与制度知识} (wuhan);
    \draw[timeline, dashed] (shanghai) -- node[right,note]
      {沿海航运与通信（示意）} (fuzhou);
    \draw[timeline, dashed] (wuhan) -- node[below,sloped,note]
      {省际学校、军队与商路网络} (guangzhou);
    \draw[source, dashed] (yunnan) -- node[below,sloped,note]
      {边境战争、交涉与地方行政} (guangzhou);

    \node[draw=source, fill=paper, rounded corners=3pt, align=center,
      text=source, font=\small] at (6.0,-3.08)
      {设省、设局、建校或铺线是制度节点，\\不等于各省同时采用或社会影响相同};
  \end{tikzpicture}
  \caption{晚清的边政、新政与省际改革网络（1878--1908，示意）。图把新疆设省、东北体制调整、北洋军政、江南工商与报刊、闽台工程、湖北新军学堂、两广商政和云南边政等可追踪节点置于同一关系图中，而不将其归结为从北京单向扩散的“现代化”。依据《清德宗实录》《清宣统政纪》、军机处和总理衙门档、《筹办夷务始末》、各省奏折与经费档、学堂／新军／铁路／电报机构记录、海关档和中文外文报刊；边疆部分并读俄、法、日及地方材料。位置与连线为约略示意，未绘行省、租借地或势力范围边界；设省、设局、建校、编练和铺线只能证明制度动作或特定节点，不能推出全省覆盖、实际控制或社会接受。}
  \label{fig:provincial-reform-networks}
\end{figure}


\section{1875—1889：边疆、灾荒与海防的不同考题}

光绪元年，皇帝即位而两宫太后仍握有关键影响，云南回民战争也在惨烈收束。次年起华北、西北的丁戊奇荒把粮价、迁徙和赈济推到地方官、绅士、商人和传教团体面前；官办、绅办与宗教网络并存，不能化为一个完整有效或全然失灵的“救灾体系”。1878年左宗棠军收复新疆大部，朝廷随即讨论设省、军政与财政：收复城镇与建立可持续的行政、供给和税收，是相连但不同的工序。

1881年伊犁条约的文本、1883—1885年中法战争中的马江与台湾战场、1885年台湾建省，分别让边界、海防与岛屿行政变成需要不断拨款和执行的事务。福建水师的惨败不能简化为“有船无用”，因为舰队训练、港口、指挥、战时情报和财政来源都在其中；台湾建省后的铁路、电报、矿务与行政计划也不等于各地居民已能平等使用其资源。清廷在这些事件中增加了机构和项目，仍无法摆脱区域经费与地方合作的限制。

西北与海岛的难题虽都被放进“自强”一类语言，实际所问并不相同。新疆在军事接收后，需要把河西走廊的饷道、绿洲城镇的市场和水利、驻防及新设州县接续起来；伊犁河谷还必须在与俄国交涉的边界、贸易和迁徙条件下处理。1881年\eventref{EQ-1881-SAINT-PETERSBURG-TREATY}{《中俄伊犁条约》}确定的是谈判双方所同意的文本安排，既不能划出居民已经接受的精确边线，也不能替代地方材料对税役、越境和日常贸易的说明。设省讨论因而不是战事的尾注：它提出了由谁筹饷、何处设官、哪些既有关系可以被纳入的问题；答案在不同绿洲、牧地和城镇并不一致。

台湾的情形也不能用同一把尺子衡量。建省把港口、防务和行政提升为朝廷明确经营的对象，刘铭传随后推动的工程试图把基隆、淡水、台北与中南部的交通和讯息接得更紧。然而工程奏请、拨款、征地、施工和实际使用属于不同环节。港口商人、平原移民聚落与内山原住民社群同铁路、矿务和抚垦政策相遇的方式不可能由一份省制文件概括。清廷文书可以让我们看见命令怎样沿部院和督抚系统下达；地方档案、契约与工程记录才可能在具体地点追问土地、劳力和设施如何被分配。

这些边地经验也改变了财政的含义。军费不是抽象的“国家投入”，而是银粮、转运、关卡、商路和经手者能否在远距离上不断接续的问题；海防同样依赖煤炭、船坞、港口和讯息，而不只依赖购得的军舰。朝廷能够命名一项工程或划出一笔款项，却仍须同省级留存、商人信用和地方征发相互折冲。正是在这种不能由中央单独解决的资源关系中，后来关于借款、铁路和路权的争执逐步获得了政治重量。

\section{1894—1901：战败不是一瞬，改革也不是同一条路}

1894年朝鲜的东学农民战争使清日派兵和区域竞争升级为甲午战争。黄海战事、辽东与海防的失利，最终导向\eventanchor{EQ-1895-SHIMONOSEKI-TREATY}{《马关条约》（1895）}：割让台湾、澎湖，承认朝鲜独立，并承受赔款和通商条件。条约能确认签约各方的承诺；台湾的抵抗、日军登陆与地方社会的遭遇仍必须从当地与日方材料逐段追踪，不能从一纸文本推得。

战败后的借款、铁路、租借地和新式学堂讨论，把国家财政同跨国资本、港口网络和省级税源拉得更紧。1898年百日维新与政变表明，诏令中的教育、官制和经济语言要通过人手、经费和地方程序才可能落地。1900年义和团运动、北京失陷和东南互保又显示，宗教冲突、地方保护、外交武力和督抚的选择无法纳入一条中央命令。1901年\eventanchor{EQ-1901-BOXER-PROTOCOL}{《辛丑条约》}把赔款、驻军和外交条件写入新一轮约束，也迫使清廷重新思考机构与财政。

赔款与借款首先是如何把未来收入变成眼前可动用款项的问题。中央要面对对外付款和信用，省级官员则要在本地税源、商贸秩序和既有开支之间安排解款；同一项财政安排因而会在海关口岸、厘金关卡和省城呈现为不同的压力。不能仅从条约列出的义务推定某一地区承担了什么，也不宜把一份预算当作款项已经征足、解足和用足的证明。不过，财政渠道被重新估量这一事实，使铁路不再只是关于速度的技术提案，而成为谁能筹资、谁可管理、收益流向何处的制度问题。

早在1892年，\eventanchor{EQ-1892-RAILWAY-DEBATE}{铁路修筑之争}已把官员、商人和士绅带到资金、主权与线路的同一张议程上。赞成者未必赞成同一种筹款方式，反对者也未必拒绝交通本身；报刊和奏议所记录的是各方公开提出的方案，而不是一幅整齐的民意图。甲午战后，这个问题与东北的外部约束尤为紧密：1896年的中俄密约及东清铁路安排，为清廷提供了对日后外交的一条路径，却也让铁路沿线的行政、治安和主权不能仅按北京的意愿铺陈。铁路既联结城市与港口，也把谈判桌上的权利、公司账簿里的资本和沿线社会的土地与劳力放在同一条线上。

同一时期的新式教育不是战争的自然产物，而是多种人群在不同地点作出的制度选择。新学堂、译书和留学的倡议可以从甲午后的奏折、文集和报刊中辨认出来；它们改变了何种知识可被当作仕途、技术或公共事务的准备，却没有立即取代私塾、书院和家庭资助。家庭能否承担学费、旅费和失去劳力的代价，地方是否有教师、校舍和稳定经费，都会改变一纸章程在县城或口岸的意义。新教育所形成的是竞争资格的重新排列，而不是一条人人同步进入的道路。

百日维新的紧张正在于，朝廷试图把这些彼此纠缠的问题迅速改写为诏令。它留下了关于官制、学堂、实业和军政的语言，也留下了短暂政治窗口内任命、反对与政变的痕迹；不能据此把未及试办的方案写成既成制度。1900年东南督抚同外国方面协商维持部分地区秩序，则把另一层现实摆到台前：督抚维护商路和城市安定的选择不等于各省脱离清廷，却显示中央战争命令、地方财政和对外风险已不再能由同一套优先次序处理。改革由此成为省级政治的一部分，而非北京单向输出的命令。

\section{1902—1908：把新制度放进旧社会}

1902年以后，新政在学制、军制、商政、警政和地方行政中陆续启动；总理衙门改为外务部，制度名称改变了处理事务的接口，却不能保证省县已有专业人员和稳定预算。1905年\eventanchor{EQ-1905-ABOLITION-OF-CIVIL-SERVICE-EXAM}{废除科举}，重写了读书、家庭投资和入仕的资格。私塾、宗族资源和地域差异并未同时消失，新式教育因而既扩大了新路径，也加深了进入路径的竞争。

预备立宪、新军编练、修律和东三省改制在1906—1907年并进。东北又在日俄战争中成为战场，清廷名义上的中立无法替当地居民隔开军事、铁路和行政压力。1908年光绪帝与慈禧太后相继去世，溥仪继位；《钦定宪法大纲》所规定的未来时间表，至此必须由一个幼帝、摄政集团、省级财政和不断扩张的公共政治共同承受。

新政并非从北京的一个部门直线抵达基层。外务部、学部、商部以及后来的邮传部改变了事务被命名、呈报和分派的门径；省县却要在旧有衙门、地方筹款和新旧人员之间找到能运转的做法。学制的发布、学校的开办和课程能否持续是三件事。废科举特别清楚地显示了这一层次：它结束的是国家考试制度，并没有立即消除以经典教育为中心的训练、地方名望或家庭对仕进的期待。对一部分青年而言，新学校、留学和新式职业成为可见选择；对另一些家庭，最先感到的也许是原有资格路径被关闭而替代资源仍然稀少。

军制改革同样要放在地方的尺度上读。1906年\eventanchor{EQ-1906-NEW-ARMY-REFORM}{新军编练}与军事学堂扩展，试图以新式训练、军官教育和编制来处理旧军制暴露出的弱点。部令能够规定操练和名目，却不能使装备、饷项、教官和军纪在每一省同时具备，也不能预先规定军官、士兵会在日后的政治危机中作何选择。新军之所以成为重要的政治空间，不是因为它天然通向革命，而是因为省级财政、学校网络、同乡与职务关系把军事组织嵌入了新的公共生活；不同驻地的经历必须分别追问。

\eventanchor{EQ-1906-CONSTITUTIONAL-PROMISE}{预备立宪的宣布（1906）}为改革添上另一重时间表。清廷需要显示制度可以更新，同时又试图把权力调整置于由中央控制的顺序之内。1908年的《钦定宪法大纲》可以证明朝廷公布了怎样的框架，不能证明代表机构已经拥有同等的权力，更不能把“预备”的日期直接写成社会信任的日期。1907年东三省改制尤其显出这种距离：行省体制旨在加强边防和行政经营，但日俄战争后铁路和港口周边仍受外国势力与既有安排牵制。制度重组和有效控制可以并存，也可以彼此抵触。

边地使新政的限度更清晰。日俄战争发生在清廷宣称中立的东北，铁路、港口与村镇却无法置身战场之外；来自俄日的军政文书也都首先服务各自的战争与外交。西北、蒙古与青藏高原的治理同样不应被东北的省制经验代替：交通条件、宗教机构、跨境贸易和行政语言不同，中央改革抵达这些空间的速度与结果也不同。所谓“近代化”的统一名称，遮蔽的正是这些不同的约束。

\section{1909—1911：咨询政治与路权的结}

1909年\eventref{EQ-1909-PROVINCIAL-ASSEMBLIES}{各省咨议局开办}，使原先主要在衙门、书院、商会或报刊中展开的论政，获得了一种可见的制度场域。它并不等于普遍代表：资格、选举和权限都有边界，地方精英也不持同一种立场。但预算、教育、铁路和地方自治因而有了可以公开陈述、组织联署并与督抚交涉的场所。1910年资政院开院，显示朝廷愿意安排中央层面的咨询程序；咨询权与实际决策权之间的距离，反而使时间表、责任和财政成为更尖锐的争论。

路权把这些积累的矛盾集中在一项人人看得见、却无人能单独完成的工程上。各省商绅、公司、官员和投资者所说的“收回”并非一个利益完全相同的群体；有人着眼于地方筹资与经营，有人担忧借款条件，有人要求中央承担干线责任。1911年\eventref{EQ-1911-RAILWAY-NATIONALIZATION}{铁路国有化的宣布}试图以中央名义整合干线并筹措外债。公告所表达的是政策方向，随后各地的抗议、谈判与处置则是另一段历史，必须回到省档、公司记录和报刊的具体地点与日期。把二者混为一事，既会低估中央财政的困境，也会抹平地方股东、学生、商人和官员在危机中的不同计算。

这条铁路争论也让晚清改革的制度得失显露出来。新学校、新军、咨议局与邮电铁路网络扩大了受教育、受训、联络和公开议论的渠道；它们同时需要稳定的经费、可信的程序和能在中央与省之间被承认的责任分配。当财政承诺、主权焦虑和地方投资在同一条线路上相撞时，改革创造的组织能力便可能转为质询改革本身的能力。1911年的政治断裂并非由一项政策机械造成，却不能脱离此前数十年中边防、赔款、铁路、省政与新式组织彼此牵动的历史。

\begin{debate}[战败或诏令，能否解释晚清改革的成败]
甲午战败、维新政变与新政章程都留下清楚的政治节点，但不能由此推出一条全国同步的“改革成功”或“改革失败”曲线。朝廷档案、条约和预算分别记录命令、承诺与拨款；学校、军队、警政和地方财政如何运行，仍须以各地材料逐项检验。
\end{debate}

\paragraph{材料位置。} 《清德宗实录》、宫中朱批奏折、内阁大库题本与预算材料能追踪朝廷的任命、章程与拨款；条约、中外照会和外交档案可限定国际交涉的文本。灾荒、战争和学校的实际经验还需要地方档案、报刊、学校名册、公司记录和家庭材料。统计、伤亡、民意与“改革成败”均不能由诏令、报纸社论或单一外交报告直接推出。

\section{战败之后：改革为何成为地方政治}

1894年甲午战争爆发，1895年《马关条约》把赔款、通商与台湾割让推到清廷必须处理的财政和主权危机中。台湾的抵抗与日军登陆之间存在一段具体而暴力的地方过程，不能由条约签署取代。赔款、借款、租借地和铁路议题随即把中央财政同全球资本和港口网络接合；它们既扩大了技术、信息和商业的渠道，也加重了对税源、土地和地方合作的争夺。

1898年变法与政变显示改革并没有单一阵营或同一速度。诏令可以建立新的官制、教育和经济语言，不能证明省县已有经费、人员和执行程序。义和团运动、1900年北京陷落和1901年条约进一步说明宗教冲突、地方保护、外交武力与国家决策不能被一个标签解释。官府的“拳匪”、教会的“迫害”、列强的“保护”都是行动中的语言，而不是完整因果。

新政把学校、军队、警察、法律、商政和地方自治带入同一轮制度试验。1905年废科举重写了读书和入仕的资格，但私塾、家庭资源与地域不平等不会因此消失；新军和新式学堂也提供政治组织的空间，却不预定每位学生或士兵会走向革命。改革既增添国家能力，也暴露它需要省级财政、专业人员和社会信任才能延续。

\subsection{材料与争议}

上谕、章程、预算、报刊与外交照会应分别作为制度宣布、资金安排、公共争论和谈判立场来读。学校名册、警察档案、公司记录、地方诉讼和家庭材料只能局部检验实践。将改革描述为必然的现代化，或将其描述为全然失败，都会忽略参与扩大与门槛并存的事实。

\section{收复、设省与海岛：边疆被重新安排}

一八七〇年代末，清廷同时背着几条很不相同、却彼此牵连的路：灾区的粮路从仍能通行的河道和商号向外伸展；左宗棠部的军需要越过甘肃与河西；伊犁的照会在北京、圣彼得堡和河谷之间往返；福建与台湾的港口又必须计算煤炭、船只和讯息抵达的先后。它们不能合成一部从北京向边缘展开的行政史。每一条路都要靠不同的人、货物和中介维系，也都把风险留在沿线。

战争的胜负因而只是重新安排的开端。军队进入绿洲，并不立即解决水利、市场与赋役；建省写入成案，并不使山地、港口和平原在同一日受同一种治理；条约的中外文本也不能代替关卡、牧地与市镇中实际发生的往来。晚清的边疆政治首先是一种昂贵的连接工程：朝廷试图把军令、银粮、信息和行政名义接起来，地方的运输者、商人、居民与社区则在征购、借贷、迁徙、贸易和避险中承受其不均衡的后果。

\subsection{饥荒中的粮路：救济不是一张覆盖全国的网}

\eventanchor{EQ-1876-NORTH-CHINA-FAMINE}{华北、西北饥荒（1876）}发生时，旱情并不是唯一的压力。粮价、仓储、道路和地方可动用的信用共同决定一袋粮能否离开产地、经过关卡或市镇、到达一个正在断炊的村庄。对离乡者而言，赈粮只是几个有限选择中的一个；投亲、借贷、出售物品、结伴迁徙或留在原处等待，都会随家庭劳力、可走道路和地方关系而改变。把灾区写成静待命令抵达的平面，便看不见这些被迫作出的判断。

到一八七七年，\eventanchor{EQ-1877-FAMINE-RELIEF-NETWORKS}{灾区赈济网络}把清廷的赈务、地方绅商的募捐、商路的转运和传教团体的救助接入同一场危机。它们有时在同一市镇相遇，有时只抵达各自能够取得消息、粮源和保护的地点。募捐簿上的一笔银两、奏报中的一批粮食，首先记录的是某项筹措或交付；它不能自动说明谁领到粮、一个家庭能维持多久，或邻近村庄是否被遗漏。灾民留下的自述本已稀少，迁徙者、儿童和无力立券记账者尤其难以在保存下来的文书中发声。

救济还改变了信息的价值。能够向县城、会馆或通商口岸传出消息的地点，更容易进入募捐者、报刊和外来救助者的视野；远离干道的村落未必因此没有互助，却较少留下可被后人追踪的记录。地方绅士的名册会突出捐输者，传教网络的报告会说明本机构的行动，督抚的奏报则要让上级看见可处理的灾情。把这些材料并置，不是要从中凑出一个无误差的死亡数，而是要辨认粮食、信用和消息曾在哪里相互接上，又在哪里中断。

这场饥荒与随后西征并非同一条因果线，却在同一个财政世界中争夺可用的粮、银与运输能力。省级官员、商人和地方组织既要应付本地的救济与治安，也要面对跨省筹饷和采购带来的价格、劳力与道路压力。清廷能够调动资源，并不意味着成本在各地平均分摊；正是在这种不均衡中，军事收复和灾后生计同时变得可能，也同时留下新的债务与脆弱性。

\subsection{河西饷道、绿洲接收与伊犁谈判}

\begin{event}{\eventanchor{EQ-1878-Zuo-ZONGTANG-XINJIANG}{收复新疆大部：饷道如何决定军队能够走多远}}{光绪四年（一八七八年）}{左宗棠部、清廷与新疆各地行动者 · 甘肃、河西走廊、吐鲁番、喀什噶尔及北路牧地}{清军收复大部与阿古柏政权崩解可核；战事日期、经费、人员与各地接收的范围仍须逐件复核}
左宗棠部继续西进时，前线的胜负始终系在后方能否把银粮、牲畜、道路修缮和递送文书连成一条不断裂的线。甘肃与河西不是地图上通往新疆的一段空白：地方官要筹解，商人和运输者要估量货价、牲力与治安，沿线居民则要面对征购、劳役和驻军需求。军队能抵达某座城，说明这条链在某一时刻仍可维持；它不表示每一笔饷项都如数到位，更不表示沿路社会可以不受代价地提供它。

阿古柏政权崩解后，吐鲁番、喀什噶尔、伊犁与北路牧地也没有进入同一个整齐的“战后”。军队到达、城镇市场重开、州县与驻防设置、水利和土地关系如何继续，各有自己的节奏。清廷由此重新获得安排行政的机会，却必须先决定何处优先供给、怎样与既有的地方关系打交道、哪些道路足以维持驻军。设省不是凯旋之后附加的一枚印章，而是要把临时军需转成日常征解、司法和交通时才显露出来的难题。

清方军报、军机处文书和地方呈报可以勾勒调兵、筹饷与战况的次序；它们不能替绿洲与牧地中的人说明服役、贸易或迁居如何改变。俄文材料、维吾尔文和蒙古文文书以及地方档案保留了另一组观察位置，但保存不全、地名与纪年转换亦有难处。因而，“收复”在这里应当先读作清军所确认的军事过程，而不是把不同地点此后的服从、负担和安全感预先写成结论。
\end{event}

设省后的新行政也要在旧有的生活空间中落脚。州县、驻防、屯垦和军台可以在文书上各有职责，实际运作却仍要依赖掌握水源、土地、市场和道路信息的人。一个税名的设立、一次驻军的调动或一段军台的整修，可能为某处带来新的交易机会，也可能增加征发与纠纷；不能把行政设置直接等同于居民已经获得相同的司法、教育或赋役经验。新省所要处理的，正是把战时的远距离供给改造成可以反复征解、传递与裁断的日常程序，而这种转换从来不在地图上同时完成。

军事压力转入外交，仍没有让这条饷道失去作用。清俄围绕伊犁归还和边界条件持续交涉的\eventanchor{EQ-1880-ILI-NEGOTIATIONS}{一八八〇年谈判}，发生在北京与圣彼得堡的照会网络中，也牵动伊犁河谷的牧地、城镇和跨境商路。清廷需要把西征取得的军政位置转换为可被承认的条件；俄国方面则将边境安全、商业与帝国利益写入自己的交涉语言。河谷中的商队、迁徙家庭和地方中介却未必依照两座首都的节奏行动：一项条文若要成为关卡的惯例，还要经过翻译、传达、征税和实际通行。

\eventanchor{EQ-1881-SAINT-PETERSBURG-TREATY}{《中俄伊犁条约》（1881）}使伊犁大部归还、赔款与贸易条件获得了条约形式。签署、批准和中俄文本能够确定双方作出了哪些承诺；它们并不能画出居民已经接受的精确边线，也不能把签约日变成地方接收完成日。对清廷而言，条约减轻了继续以战场方式处理伊犁的压力，却把跨境贸易、赔款和边界管理一并带进财政与行政的日常账目。对河谷社会而言，哪些货物能够过关、谁要承担税役、牧地和亲属网络如何被新手续切分，仍须从地方与多语材料中逐处追问。

条约中的赔款与通商条件也提醒人们，外交并非发生在军需之外。赔付需要筹措和解送，贸易条款需要在口岸、关卡和账册中被解释，边界管理则要不断处理来往的人、牲畜和货物。清廷和俄国的照会能清楚呈现各自要坚持的语言，却较难记录一笔交易因新手续拖延多久、一户迁徙家庭怎样重排亲属关系。这里的证据限制不只是史料学的附注；它改变了对“定界”一词的读法，使边界成为持续被执行、规避和协商的过程。

\subsection{从北圻到台湾：海防是一组港口、道路和账目的选择}

\begin{event}{\eventanchor{EQ-1883-SINO-FRENCH-WAR-BEGIN}{北圻交涉变成战争}}{光绪九年（一八八三年）}{清、法国、越南及两广边路行动者 · 北圻、红河、两广与北京}{冲突升级为战争可核；开战责任、地方控制与兵力规模不能由任一方档案裁定}
北圻的冲突升级时，红河水路、城镇、两广边路与海防港口同时决定军需和消息能够怎样移动。清廷面对的并不只是关于宗藩关系的外交措辞：边将要判断何处可驻、何处可运粮，越南朝廷和地方行动者也在法国军政推进与本地关系之间作出选择。法国拥有海军优势，并不等于它无需把炮火、河口、登陆、补给和城镇维持接在一起。各方都在有限讯息下行动，战争因而不是两份声明的自动结果。

北圻的战局也把福建、广东和台湾卷入同一场资源竞争。两广边路所需的转运与驻军，港口防务所需的煤炭、修理和船员，都会占用本可用于别处的银钱与注意力。北京发出的调兵令能够显示它试图如何排列优先次序，却不能替每一段边路说明翻译、商路、雨季和地方避险怎样改变了命令的落点。把越南只写成清法竞争的舞台，或把沿海社区只写成海军的背景，都会抹去这些相互牵制的时间。
\end{event}

\eventanchor{EQ-1884-FUZHOU-NAVAL-BATTLE}{马江海战（1884）}使福州船政、港湾防御和法国舰队在短时间内正面相遇。船只并非脱离训练、指挥、煤炭、港口和情报而独立发挥作用；一支舰队即使在纸面上拥有船只，也要能让人员、弹药、修理和讯息同时抵达。战后，海防经费的去向不再只是海图上的问题：船坞附近的工匠、运输煤粮的船户、依港口市场谋生者以及承担加派的腹地，都要在战争留下的支出中寻找位置。中法军报、船政记录和地方材料对于舰只损失与伤亡并不总能相互吻合，尤其不能用一份战报的数目概括港湾内外的全部经验。

台湾因此不是这场战争的附带地名。封锁与基隆、沪尾一带的防御迫使决策者把港口、煤源、粮运、电报与炮台放在同一张预算中计算。战后，\eventanchor{EQ-1885-TAIWAN-PROVINCE}{台湾建省（1885）}把军防、行政和工程纳入新的制度计划。刘铭传推动铁路、电报、矿务与抚垦时，计划、拨款、征地、施工和日后使用是彼此不同的步骤：淡水、基隆和台南的港市，平原移民聚落和内山各社群与这些步骤的关系不会相同。省名可以表明朝廷重新排列行政与防务的意图，不能压平岛内因土地、劳役、道路和武力接触而出现的差异。

工程尤其使行政的尺度变得可见。一段铁路或电报线要先取得资金、材料、测量和施工的条件，才可能改变何处接到消息、何处运输较快；矿务的设想也要经过土地、劳力和市场才会产生可见后果。地方档案与契约能够让我们在一些地点看见征地或劳役的痕迹，不能把这些片段推广成全岛统一的经验。内山社群的声音在留存材料中常被军政与开发文书遮蔽，正是“经营台湾”不能被写成一项没有摩擦的行政成就的原因。

\subsection{朝鲜危机、甲午战场与海峡的后果}

\begin{event}{\eventanchor{EQ-1894-FIRST-SINO-JAPANESE-WAR}{甲午开战：朝鲜、黄海与辽东的补给压力}}{光绪二十年（一八九四年）}{清、日本与朝鲜行动者 · 朝鲜半岛、黄海、辽东及沿岸港口}{开战及区域扩展可核；军事规模、责任和不同社会的经历不能由一方战报判定}
朝鲜的政治行动、驻军与外交交涉先改变了危机的条件，不能把朝鲜缩成清日竞争的一块背景板。随后黄海的舰队、辽东的道路和沿海港口把煤炭、粮食、情报与伤员运输拖入不同战场。清廷与日本都试图以调兵、筹饷和文书把分散的资源接向前线；船只是否可用、港口能否补给、道路是否容得下辎重，却不由北京或东京的一道命令决定。

这场战争的直接后果是区域秩序的军事化升级，较长的代价则落在海防财政、边地居民与跨境贸易网络上。舰队兵额、战果数字和将领责任，在中日朝文及外文记录中各有用途与计算方式；它们能够帮助确认战争怎样扩展，却不能使研究者替不同地点的居民说出相同的感受。把朝鲜、黄海和辽东并置，所要看见的正是同一场开战如何在不同的补给条件中变形。
\end{event}

一八九五年的\eventref{EQ-1895-SHIMONOSEKI-TREATY}{《马关条约》}规定了割让台湾、澎湖以及赔款和通商等条件，但海峡两岸的局面没有随着条约文字停止。清方撤离、日本登陆、港口交通与地方组织之间留下了短促而充满不确定的空隙。有人试图组织抵抗，有人要保护家人、货物和村落，有人被迫随交通线或战事移动；将这些不同处境都称作一个一致的“地方反应”，反而会再次遮住岛内的差异。

\begin{event}{\eventanchor{EQ-1895-TAIWAN-REPUBLIC-RESISTANCE}{台湾民主国与抵抗日军登陆}}{光绪二十一年（一八九五年）}{台湾地方行动者、清与日本 · 港口、道路、城镇与乡村}{建国名义及抵抗发生可核；组织范围、地方选择和战事过程须以中日及地方材料互校}
地方力量建立“台湾民主国”并抵抗日军登陆，表明签约后的主权转换仍须经过港口、道路、城镇与补给线。名义上的建国和实际能够动员、守住或离开哪些地点，并不是同一回事。清廷、日本方面和地方记录使用的政治语言也各不相同；它们可以标出组织、军队与行政所声称的行动，却不能使任何一方的措辞代表全岛居民的判断。

这一段抵抗不是条约自动产生的反应。铁路、电报和港口设施曾被设想为加强岛内联系的工具，战时却同样成为移动、传讯与争夺的条件；它们的作用取决于何处修成、何处可用以及谁能接近。割让后的创伤因此不能从一纸文本读尽，也不能只由登陆战决定。土地契约、地方档案、工程记录与日文军政材料能够让研究者在若干地点追踪征地、劳役、航线和武力接触，但原住民社群与普通居民的自述保存尤其有限，留在材料之外的经验必须被明确保留。
\end{event}

\subsection{看见连接，也看见断裂}

晚清的边疆与海岛政策并非没有效力：军队重新进入新疆，伊犁谈判形成条约，台湾的行政、通信和工程获得新的组织方式，港口与边路也被纳入更密集的军事筹算。这些安排解决了清廷在特定时刻的通行、供给与防务问题；它们的副作用同样具体。长距离饷运把征购和信用压力推向沿线，设省把土地、劳役与行政重新分配，条约和战争又使边境市镇、商队和港口社会必须适应新的手续与风险。

可以保存下来的文书多半来自正在调度、请款、谈判或归责的机构。它们善于留下命令怎样起草、银粮怎样呈报、条约怎样措辞，也较少留下未被登记的迁徙、拒绝、协商和家庭离散。地方档案、契约、价格记录、港口材料以及俄文、维吾尔文、蒙古文、越南文、日文和其他外文材料能够补足若干地点，却不构成一张没有空白的全景。正因如此，收复、建省、海战和割让都应被读成仍在进行的区域过程：国家的线路向外延伸，地方社会却从未只是线路末端的被动注脚。

